\section{The EVM CLOB Facade}
\label{sec:precompile}

Solidity contracts and EVM tooling expect a pool-manager interface. The D-Chain exposes
one through a native precompile at address
\texttt{0x0000000000000000000000000000000000009010}. The precompile is a facade: it
presents a Uniswap-V4-\texttt{PoolManager}-style surface but routes every call into the
unified D-Chain. It is not an embedded AMM and it holds no canonical state.

\subsection{A CLOB behind a PoolManager surface}

The crucial correction to a common misreading: \texttt{0x9010} is a \emph{central limit
order book}, not a constant-product AMM. The V4-style method names map onto CLOB
operations:

\begin{table}[H]
\centering
\renewcommand{\arraystretch}{1.25}
\begin{tabular}{p{4.0cm} p{8.0cm}}
\toprule
\textbf{V4-style method} & \textbf{CLOB meaning} \\
\midrule
\texttt{initialize} & Create the market / order book. \\
\texttt{modifyLiquidity} & Place or cancel a resting limit order. ``Liquidity'' is
resting orders, not curve reserves. \\
\texttt{swap} & Submit a marketable order against the book; the return is the
\texttt{BalanceDelta} from the resulting fills. \\
\bottomrule
\end{tabular}
\caption{The facade maps Uniswap-V4 \texttt{PoolManager} methods onto CLOB operations.
The precompile translates Solidity calls into D-Chain order operations over the frozen
ZAP wire format; it executes no matching itself.}
\label{tab:facade}
\end{table}

The precompile speaks the frozen ZAP wire format (\texttt{clob\_place},
\texttt{clob\_cancel}, \texttt{clob\_submit}) and forwards to the atomic proxy, which
relays to the D-Chain. It is a read-through adapter: it holds no canonical pools, and DEX
state lives only on the D-Chain.

\subsection{Inert by default}

The precompile ships in the public EVM but defaults to an inert backend that refuses every
call until a venue configures a D-Chain endpoint. The endpoint is opt-in configuration; an
empty endpoint leaves the DEX disabled. Only when a venue sets the endpoint does the facade
bind to the live ZAP engine and route calls to the D-Chain. This is the property that lets
the precompile exist in the open-source node without dragging in the proprietary engine or
forcing any validator to run a DEX.

\subsection{Relationship to the V2/V3 AMMs}

The V2 and V3 contracts are ordinary Solidity constant-product AMMs, self-contained and
separate from the precompile. They are not the facade and do not route into the D-Chain.
Native LUX is \texttt{address(0)} as the canonical key; wrapped LUX, where used, is a
distinct asset. Keeping the AMMs orthogonal to the CLOB facade preserves the
one-home-per-concern discipline: the curve-based path and the order-book path do not share
state or implementation.

\subsection{Why a precompile rather than a contract}

Running the facade as compiled node code rather than interpreted EVM bytecode removes
opcode-dispatch and ABI-decode overhead from the hot path, and---because matching does not
happen in the precompile at all---removes per-validator re-execution of matching from the
EVM. The precompile's only job is to translate and forward; matching and state are the
D-Chain's job. This separation is what allows the matching engine to be a GPU pipeline
(Section~\ref{sec:gpu-matching}) instead of being bounded by the slowest validator's serial
EVM.
